% Options for packages loaded elsewhere
% Options for packages loaded elsewhere
\PassOptionsToPackage{unicode}{hyperref}
\PassOptionsToPackage{hyphens}{url}
\PassOptionsToPackage{dvipsnames,svgnames,x11names}{xcolor}
%
\documentclass[
  letterpaper,
  DIV=11,
  numbers=noendperiod]{scrartcl}
\usepackage{xcolor}
\usepackage{amsmath,amssymb}
\setcounter{secnumdepth}{-\maxdimen} % remove section numbering
\usepackage{iftex}
\ifPDFTeX
  \usepackage[T1]{fontenc}
  \usepackage[utf8]{inputenc}
  \usepackage{textcomp} % provide euro and other symbols
\else % if luatex or xetex
  \usepackage{unicode-math} % this also loads fontspec
  \defaultfontfeatures{Scale=MatchLowercase}
  \defaultfontfeatures[\rmfamily]{Ligatures=TeX,Scale=1}
\fi
\usepackage{lmodern}
\ifPDFTeX\else
  % xetex/luatex font selection
\fi
% Use upquote if available, for straight quotes in verbatim environments
\IfFileExists{upquote.sty}{\usepackage{upquote}}{}
\IfFileExists{microtype.sty}{% use microtype if available
  \usepackage[]{microtype}
  \UseMicrotypeSet[protrusion]{basicmath} % disable protrusion for tt fonts
}{}
\makeatletter
\@ifundefined{KOMAClassName}{% if non-KOMA class
  \IfFileExists{parskip.sty}{%
    \usepackage{parskip}
  }{% else
    \setlength{\parindent}{0pt}
    \setlength{\parskip}{6pt plus 2pt minus 1pt}}
}{% if KOMA class
  \KOMAoptions{parskip=half}}
\makeatother
% Make \paragraph and \subparagraph free-standing
\makeatletter
\ifx\paragraph\undefined\else
  \let\oldparagraph\paragraph
  \renewcommand{\paragraph}{
    \@ifstar
      \xxxParagraphStar
      \xxxParagraphNoStar
  }
  \newcommand{\xxxParagraphStar}[1]{\oldparagraph*{#1}\mbox{}}
  \newcommand{\xxxParagraphNoStar}[1]{\oldparagraph{#1}\mbox{}}
\fi
\ifx\subparagraph\undefined\else
  \let\oldsubparagraph\subparagraph
  \renewcommand{\subparagraph}{
    \@ifstar
      \xxxSubParagraphStar
      \xxxSubParagraphNoStar
  }
  \newcommand{\xxxSubParagraphStar}[1]{\oldsubparagraph*{#1}\mbox{}}
  \newcommand{\xxxSubParagraphNoStar}[1]{\oldsubparagraph{#1}\mbox{}}
\fi
\makeatother


\usepackage{longtable,booktabs,array}
\usepackage{calc} % for calculating minipage widths
% Correct order of tables after \paragraph or \subparagraph
\usepackage{etoolbox}
\makeatletter
\patchcmd\longtable{\par}{\if@noskipsec\mbox{}\fi\par}{}{}
\makeatother
% Allow footnotes in longtable head/foot
\IfFileExists{footnotehyper.sty}{\usepackage{footnotehyper}}{\usepackage{footnote}}
\makesavenoteenv{longtable}
\usepackage{graphicx}
\makeatletter
\newsavebox\pandoc@box
\newcommand*\pandocbounded[1]{% scales image to fit in text height/width
  \sbox\pandoc@box{#1}%
  \Gscale@div\@tempa{\textheight}{\dimexpr\ht\pandoc@box+\dp\pandoc@box\relax}%
  \Gscale@div\@tempb{\linewidth}{\wd\pandoc@box}%
  \ifdim\@tempb\p@<\@tempa\p@\let\@tempa\@tempb\fi% select the smaller of both
  \ifdim\@tempa\p@<\p@\scalebox{\@tempa}{\usebox\pandoc@box}%
  \else\usebox{\pandoc@box}%
  \fi%
}
% Set default figure placement to htbp
\def\fps@figure{htbp}
\makeatother


% definitions for citeproc citations
\NewDocumentCommand\citeproctext{}{}
\NewDocumentCommand\citeproc{mm}{%
  \begingroup\def\citeproctext{#2}\cite{#1}\endgroup}
\makeatletter
 % allow citations to break across lines
 \let\@cite@ofmt\@firstofone
 % avoid brackets around text for \cite:
 \def\@biblabel#1{}
 \def\@cite#1#2{{#1\if@tempswa , #2\fi}}
\makeatother
\newlength{\cslhangindent}
\setlength{\cslhangindent}{1.5em}
\newlength{\csllabelwidth}
\setlength{\csllabelwidth}{3em}
\newenvironment{CSLReferences}[2] % #1 hanging-indent, #2 entry-spacing
 {\begin{list}{}{%
  \setlength{\itemindent}{0pt}
  \setlength{\leftmargin}{0pt}
  \setlength{\parsep}{0pt}
  % turn on hanging indent if param 1 is 1
  \ifodd #1
   \setlength{\leftmargin}{\cslhangindent}
   \setlength{\itemindent}{-1\cslhangindent}
  \fi
  % set entry spacing
  \setlength{\itemsep}{#2\baselineskip}}}
 {\end{list}}
\usepackage{calc}
\newcommand{\CSLBlock}[1]{\hfill\break\parbox[t]{\linewidth}{\strut\ignorespaces#1\strut}}
\newcommand{\CSLLeftMargin}[1]{\parbox[t]{\csllabelwidth}{\strut#1\strut}}
\newcommand{\CSLRightInline}[1]{\parbox[t]{\linewidth - \csllabelwidth}{\strut#1\strut}}
\newcommand{\CSLIndent}[1]{\hspace{\cslhangindent}#1}



\setlength{\emergencystretch}{3em} % prevent overfull lines

\providecommand{\tightlist}{%
  \setlength{\itemsep}{0pt}\setlength{\parskip}{0pt}}



 


\KOMAoption{captions}{tableheading}
\makeatletter
\@ifpackageloaded{caption}{}{\usepackage{caption}}
\AtBeginDocument{%
\ifdefined\contentsname
  \renewcommand*\contentsname{Table of contents}
\else
  \newcommand\contentsname{Table of contents}
\fi
\ifdefined\listfigurename
  \renewcommand*\listfigurename{List of Figures}
\else
  \newcommand\listfigurename{List of Figures}
\fi
\ifdefined\listtablename
  \renewcommand*\listtablename{List of Tables}
\else
  \newcommand\listtablename{List of Tables}
\fi
\ifdefined\figurename
  \renewcommand*\figurename{Figure}
\else
  \newcommand\figurename{Figure}
\fi
\ifdefined\tablename
  \renewcommand*\tablename{Table}
\else
  \newcommand\tablename{Table}
\fi
}
\@ifpackageloaded{float}{}{\usepackage{float}}
\floatstyle{ruled}
\@ifundefined{c@chapter}{\newfloat{codelisting}{h}{lop}}{\newfloat{codelisting}{h}{lop}[chapter]}
\floatname{codelisting}{Listing}
\newcommand*\listoflistings{\listof{codelisting}{List of Listings}}
\makeatother
\makeatletter
\makeatother
\makeatletter
\@ifpackageloaded{caption}{}{\usepackage{caption}}
\@ifpackageloaded{subcaption}{}{\usepackage{subcaption}}
\makeatother
\usepackage{bookmark}
\IfFileExists{xurl.sty}{\usepackage{xurl}}{} % add URL line breaks if available
\urlstyle{same}
\hypersetup{
  pdftitle={Achieved human agency or loss of choice in Agentic AI for students},
  pdfkeywords={social psychology, sociology, human agency, agentic AI,
turing trap, AI Literacy, beliefs},
  colorlinks=true,
  linkcolor={blue},
  filecolor={Maroon},
  citecolor={Blue},
  urlcolor={Blue},
  pdfcreator={LaTeX via pandoc}}


\title{Achieved human agency or loss of choice in Agentic AI for
students}
\author{Frederik J van Deventer}
\date{}
\begin{document}
\maketitle


\subsection{Introduction}\label{introduction}

Young people around the globe are growing up in a different digital
reality than their parents, where Artificial Intelligence (AI) is
playing a big part. Regardless of whether the term of ``AI natives''
bares truth, there is a new group of students that spend almost 7 hours
and 30 minutes online per day (Kemp 2025), on platforms where AI is
embedded, e.g.~in search algorithms, AI answers, or chat functions often
signified with ``AI'' or ✨. However, this does not mean that these
students have an innate skill or that they can see ``behind the
curtain'' of these technologies just as we expected them to in the
debunked ``digital natives'' narrative of the 2000s (Selwyn 2009).

Over the last decade the question whether Artificial Intelligence (AI)
itself exhibits agency has been debated (Legaspi, He, and Toyoizumi
2019; Swanepoel 2021). This distinction between human and technological
agency becomes more blurred with the emerging paradigm of Agentic AI
(AAI), which is different to previous forms of AI such as Generative AI
(GenAI), which we encounter in platforms lik ChatGPT, Claude or Co-pilot
(Naik, Naik, and Naik 2024), because AAI can operate (semi-)autonomously
and pursue complex goals with little human intervention (Acharya,
Kuppan, and Divya 2025). Both GenAI and AAI largely depend on recent
developments in Large Language Models (LLMs). With AAI multiple
instances of these GenAI-bots run in the background interfacing each
other creating seemingly autonomous and agentic systems, as a
evolutionary next step from the GenAI chat-platforms widely in use among
students (Silvennoinen, Aksovaara, and Alanko-Turunen 2025; Dai 2025).

This questions the notion of human agency (Anderson and Rainie 2023), as
systems now also seem to (1) act with (2) intention and a certain kind
of (3) reasoning, three aspects often associated with agency (Anscombe
1956; Davidson 1978). The chordal triad of Emirbayer and Mische (1998)
suggests placing agency within not just a structure (Giddens 1984) but
in what Mead (1932) calls the \emph{temporal} dimension of sociality,
using past, present and imagined futures as 3 dimensions on which agency
is performed: \emph{iterational}, \emph{practical-evaluative} and
\emph{prospective}. Most of these theories are individual and describe
the action of an agency up to a certain point: that of acting itself and
the effect it causes. Biesta and Tedder (2007) engages with the chordal
triad of agency by focussing on the notion of \emph{agency achieved} in
an \emph{``ecological understanding''} of agency (Biesta and Tedder,
n.d.); it is here that the effect of action becomes realized and
experienced.

The \emph{iterational} aspect is exposed in structure and collective
beliefs or ideology expressed in instruments such as policies, or
punitive action or attitudes of individuals. Beliefs about
``inevitability of technology'' and the cult of innovation (Winner 2018)
pervade, while at the same time mistrust and weariness tries to keeps AI
out of universities (Guest et al. 2025). Technology is never just
neutral nor is it inherently good or bad (Morrow 2014; Heyndels 2023) it
is be assistive to our human needs (Arora 2024, 32). Regulation and
policy fueled by ideology (both benefiting the few or the many) has been
the driver for the direction this advancement would take us (Johnson and
Acemoglu 2023, 57).

Students navigate these punitive faculties, policies and teacher
narratives as ecology in which they need to realize goals, their agency
changing from situation to situation. Are they seeing their volitions
realized as \emph{achieved agency}?

\subsection{Knowledge Gap}\label{knowledge-gap}

Most literature focusing on AI is not (yet) mentioning agentic AI, but
using GenAI, which is closely related as precursor to Agentic AI
paradigm that also utilizes LLMs. So a number of the studies mentioned
below are chosen from that context.

In the course of history, human development has been argued to increase
where agency and freedom increase (Prados de la Escosura 2022; Sen
1999). Agency is however a bloated term that can mean a number of things
in different domains, such as psychology where it focuses on individual
human agency (Bandura 2001) or sociology where it focuses on the
structures in which an actor operates (Giddens 1984). Agency is a term
closely related to and with overlaps to terms such as autonomy and
freedom. To be clear in this context, ``agency'' means for an actor to
be free to act true to that same actor's intentions, considering their
circumstance (structure),their beliefs (iterational aspect of Emirbayer
and Mische (1998)).

This in itself is impossible to measure cleanly, being dependent on so
many factors. But some attempts have been made to measure how people
perceive their \emph{sense of agency} (Tapal, Oren, and Eitam 2017) or
how Trust in Automation (TiA) influences there autonomy (Kohn et al.
2021). There are also different models that propose ``measures of
agency'' (Grünbaum and Christensen 2020) or how concepts like autonomy
and ability contribute to agency or measuring proxies for agency (Alkire
2008).

Agentic AI and AI goals in general, largely seem to focus on increasing
abilities for models or replacing tedious tasks performed by humans.
While research has been done about future projections of human agency in
conjunction with AI in decision making (Pew Research Center, Anderson,
and Rainie 2023) and in education (Mouta, Pinto-Llorente, and
Torrecilla-Sánchez 2025) it does not deal with how to increase human
agency for humans now. Or what the effect is of Agentic AI on agency or
Sense of Agency (SoA).

The way Agentic AI or GenAI is influencing agency in students seems not
to have been studied, apart from the influence it has on critical
thinking (Suriano et al. 2025), and why and how they use it for studying
(Dai 2025; Silvennoinen, Aksovaara, and Alanko-Turunen 2025; Suonpää,
Heikkilä, and Dimkar 2024). The Human AI - Interaction framework
suggested by Sundar (2020) incorporates agency, but mostly highlights
interesting avenues to pursue rather than diving into topics such as:
human and AI collaboration, the seeming either/or thinking of machine vs
human agency; agency is seen as a ``mediating variable'' in the
machine/human interaction (Sundar 2020).

One area of study related to increasing human agency surrounding AI is
\emph{AI Literacies}. Learning about AI systems explains how they
perform, this demystifies and increases people's insight in these
systems (Pinski and Benlian 2024) which helps to reduce anthropomorphism
(Druga and Ko 2021). Anthropomorphism is one of the things that can lead
to false perceptions and misunderstanding of AI abilities(Barrow 2024).
AI Literacies among students has not been studied in relation to neither
human agency, their sense of agency, nor other similar approaches such
as a capabilities approach.

In goal attainment and AI assisted coaching, there is a real opportunity
for an AI to participate in the process this then is suggested to
increase agency (Plotkina and Sri Ramalu 2024). A review collecting
these kinds of studies, that are influential to agency could not be
found.

Then temporal aspect of agency as suggested in the above mentioned
\emph{chordal triad of agency} seems not have been used at all towards
agentic AI or GenAI.

\subsection{Contribution to scientific
discourse}\label{contribution-to-scientific-discourse}

The threats and opportunities for the achieved agency of students are
unclear in the context of Agentic AI. First of all collecting work on
how AAI and GenAI are shaping agency needs to be reviewed. Secondly I
want to research the \emph{iterational} dimension (Emirbayer and Mische
1998, 975) by uncovering what beliefs or ideology is encapsulated in
teacher narratives, rules and policy surrounding (A)AI. Thirdly I want
to Lastly I will investigate how students enact their intentions within
their temporal-relational context of action and how their achieving of
agency feedbacks into shaping their context, by looking at how students
react to changing policies. This means looking into the
\emph{prospective} dimension of the chordal triad (Emirbayer and Mische
1998, 984), where goals or aspirations are set and then actualized
through the \emph{practical-evaluative} dimension.

\subsection{Research Questions}\label{research-questions}

In what way is \emph{agency achieved} impacted for students as a result
of changing circumstances with the rise of agentic AI?

\subsubsection{Sub-questions}\label{sub-questions}

\begin{enumerate}
\def\labelenumi{\arabic{enumi}.}
\tightlist
\item
  How can we categorize known effects of GenAI and Agentic AI on agency
  systematically?
\item
  How are prevalent beliefs disclosed in the ecology of achieved agency
  changing policy and teacher narratives because of AAI?
\item
  In what ways are students performing day-to-day tasks, by either
  conforming or subverting rules and restrictions?
\end{enumerate}

\subsection{Methodology}\label{methodology}

In the following section I will elaborate on how I will conduct the
research described in the above mentioned research questions.

\begin{longtable}[]{@{}
  >{\raggedright\arraybackslash}p{(\linewidth - 4\tabcolsep) * \real{0.0417}}
  >{\raggedright\arraybackslash}p{(\linewidth - 4\tabcolsep) * \real{0.7917}}
  >{\raggedright\arraybackslash}p{(\linewidth - 4\tabcolsep) * \real{0.1667}}@{}}
\caption{Research Questions and how they relate to
methods}\tabularnewline
\toprule\noalign{}
\begin{minipage}[b]{\linewidth}\raggedright
\end{minipage} & \begin{minipage}[b]{\linewidth}\raggedright
Question
\end{minipage} & \begin{minipage}[b]{\linewidth}\raggedright
Method
\end{minipage} \\
\midrule\noalign{}
\endfirsthead
\toprule\noalign{}
\begin{minipage}[b]{\linewidth}\raggedright
\end{minipage} & \begin{minipage}[b]{\linewidth}\raggedright
Question
\end{minipage} & \begin{minipage}[b]{\linewidth}\raggedright
Method
\end{minipage} \\
\midrule\noalign{}
\endhead
\bottomrule\noalign{}
\endlastfoot
RQ1 & How can we categorize known effects of GenAI and Agentic AI on
agency systematically? & Systematic Literature Review \\
RQ2 & How are prevalent beliefs disclosed in policy and teacher
narratives? & Critical Discourse Analyis \\
RQ3 & In what ways are students performing day-to-day tasks, by either
conforming or subverting rules and restrictions? & Mixed-Method \\
\end{longtable}

\subsubsection{RQ1: Systematic Literature
Review}\label{rq1-systematic-literature-review}

To find out what studies related to agency a systematic literature will
be conducted that searches for studies that relate to agency, such as
studies that focus on empowerment (Kong, Zhu, and Yang 2025). According
to the systems of Grünbaum and Christensen (2020) and Alkire (2008) when
studies are performed on proxies of agency, or related concepts, they
can be categorized and labeled in relation to AI or more specifically to
Agentic AI if those are available. Furthermore the way agency is
approached is important to delineat, whether it falls under sociology,
social psychology or philosophy

\subsubsection{RQ2: Critical Discourse
Analyis}\label{rq2-critical-discourse-analyis}

Through \emph{critical discourse analyis} (CDA) I will look at what kind
of language is constructing narratives, describes relations of power and
how it creates or reinforces myths (Blommaert and Bulcaen 2000).

CDA tries to find patterns in power dynamics and social effects in
discourse. In this analysis we can take the texts further and analyze
context and intertextuality to challenge social inequalities. I will
focus mainly on step two: \emph{Identify obstacles to addressing the
social wrong} of this process which consists of three stages (Fairclough
2013):

\begin{enumerate}
\def\labelenumi{\arabic{enumi}.}
\tightlist
\item
  Analyse dialectical relations
\item
  Select texts
\item
  Analyse texts interdiscursively and linguistically
\end{enumerate}

\subsubsection{RQ3: Mixed-method}\label{rq3-mixed-method}

\emph{Thematic analysis} will be used as Naeem et al. (2023) and Braun
and Clarke (2021) describe it to find common themes and topics in
personal qualitative interview with students. Thematic analysis is a
more subjective approach to data, where language is central. It requires
a \emph{coding} process that uses semantic and latent codes to
categorize statements and find \emph{themes} and \emph{subthemes} of
meaning in texts.

After the initial analysis and recognized larger themes questions will
be formulated to be able to use a more quantitative approach in the form
of a survey. practeval practeval(Practical-evaluative)
prospective(Prospective) --\textgreater{} practeval practeval
-.agency.-\textgreater{} Action

``` --\textgreater{}

\subsection*{References}\label{references}
\addcontentsline{toc}{subsection}{References}

\phantomsection\label{refs}
\begin{CSLReferences}{1}{0}
\bibitem[\citeproctext]{ref-acharyaAgentic2025}
Acharya, Deepak Bhaskar, Karthigeyan Kuppan, and B. Divya. 2025.
{``Agentic {AI}: {Autonomous} Intelligence for Complex Goals---a
Comprehensive Survey.''} \emph{IEEE Access : Practical Innovations, Open
Solutions} 13: 18912--36.
\url{https://doi.org/10.1109/ACCESS.2025.3532853}.

\bibitem[\citeproctext]{ref-alkire2008concepts}
Alkire, Sabina. 2008. {``Concepts and Measures of Agency.''}

\bibitem[\citeproctext]{ref-andersonFutureHumanAgency2023}
Anderson, Janna, and Lee Rainie. 2023. {``The {Future} of {Human
Agency}.''} \emph{Pew Research Center}.

\bibitem[\citeproctext]{ref-anscombe1956intention}
Anscombe, Gertrude Elizabeth Margaret. 1956. {``Intention.''} In
\emph{Proceedings of the Aristotelian Society}, 57:321--32. JSTOR.

\bibitem[\citeproctext]{ref-arora2024pessimism}
Arora, Payal. 2024. \emph{From Pessimism to Promise: {Lessons} from the
{Global South} on Designing Inclusive Tech}. MIT Press.

\bibitem[\citeproctext]{ref-bandura2001social}
Bandura, Albert. 2001. {``Social Cognitive Theory: {An} Agentic
Perspective.''} \emph{Annual Review of Psychology} 52 (1): 1--26.

\bibitem[\citeproctext]{ref-barrowAnthropomorphismAIHype2024}
Barrow, Nicholas. 2024. {``Anthropomorphism and {AI} Hype.''} \emph{AI
and Ethics} 4 (3): 707--11.
\url{https://doi.org/10.1007/s43681-024-00454-1}.

\bibitem[\citeproctext]{ref-biestaAgencyLearningLifecourse2007}
Biesta, Gert, and Michael Tedder. 2007. {``Agency and Learning in the
Lifecourse: {Towards} an Ecological Perspective.''} \emph{Studies in the
Education of Adults} 39 (2): 132--49.
\url{https://doi.org/10.1080/02660830.2007.11661545}.

\bibitem[\citeproctext]{ref-biestaHowAgencyPossible}
---------. n.d. {``How Is Agency Possible? {Towards} an Ecological
Understanding of Agency-as-Achievement.''}

\bibitem[\citeproctext]{ref-blommaertCriticalDiscourseAnalysis2000}
Blommaert, Jan, and Chris Bulcaen. 2000. {``Critical {Discourse
Analysis}.''} \emph{Annual Review of Anthropology} 29: 447--66.
\url{https://www.jstor.org/stable/223428}.

\bibitem[\citeproctext]{ref-braun2021thematic}
Braun, Virginia, and Victoria Clarke. 2021. {``Thematic Analysis: A
Practical Guide to Understanding and Doing.''} \emph{Thousand Oaks}.

\bibitem[\citeproctext]{ref-daiWhyStudentsUse2025}
Dai, Yun. 2025. {``Why Students Use or Not Use Generative {AI}:
{Student} Conceptions, Concerns, and Implications for Engineering
Education.''} \emph{Digital Engineering} 4 (March): 100019.
\url{https://doi.org/10.1016/j.dte.2024.100019}.

\bibitem[\citeproctext]{ref-davidsonIntending1978}
Davidson, Donald. 1978. {``Intending.''} In \emph{Philosophy of
{History} and {Action}: {Papers Presented} at the {First Jerusalem
Philosophical Encounter December} 1974}, edited by Yirmiahu Yovel,
41--60. Dordrecht: Springer Netherlands.
\url{https://doi.org/10.1007/978-94-009-9365-5_5}.

\bibitem[\citeproctext]{ref-drugaHowChildrensPerceptions2021}
Druga, Stefania, and Amy J Ko. 2021. {``How Do Children's Perceptions of
Machine Intelligence Change When Training and Coding Smart Programs?''}
In \emph{Interaction {Design} and {Children}}, 49--61. Athens Greece:
ACM. \url{https://doi.org/10.1145/3459990.3460712}.

\bibitem[\citeproctext]{ref-emirbayerWhatAgency1998}
Emirbayer, Mustafa, and Ann Mische. 1998. {``What {Is Agency}?''}
\emph{American Journal of Sociology} 103 (4): 962--1023.
\url{https://doi.org/10.1086/231294}.

\bibitem[\citeproctext]{ref-fairclough2013critical}
Fairclough, Norman. 2013. \emph{Critical Discourse Analysis: {The}
Critical Study of Language}. Routledge.

\bibitem[\citeproctext]{ref-giddens1984constitution}
Giddens, Anthony. 1984. \emph{The Constitution of Society: {Outline} of
the Theory of Structuration}. Univ of California Press.

\bibitem[\citeproctext]{ref-grunbaumMeasuresAgency2020}
Grünbaum, Thor, and Mark Schram Christensen. 2020. {``Measures of
Agency.''} \emph{Neuroscience of Consciousness} 2020 (1): niaa019.
\url{https://doi.org/10.1093/nc/niaa019}.

\bibitem[\citeproctext]{ref-guestUncriticalAdoptionAI2025}
Guest, Olivia, Marcela Suarez, Barbara Müller, Edwin van Meerkerk,
Arnoud Oude Groote Beverborg, Ronald de Haan, Andrea Reyes Elizondo, et
al. 2025. {``Against the {Uncritical Adoption} of '{AI}' {Technologies}
in {Academia}.''} Zenodo. \url{https://doi.org/10.5281/ZENODO.17065099}.

\bibitem[\citeproctext]{ref-heyndelsTechnologyNeutrality2023}
Heyndels, Sybren. 2023. {``Technology and {Neutrality}.''}
\emph{Philosophy \& Technology} 36 (4): 75.
\url{https://doi.org/10.1007/s13347-023-00672-1}.

\bibitem[\citeproctext]{ref-johnsonPowerProgressOur2023}
Johnson, Simon, and Daron Acemoglu. 2023. \emph{Power and {Progress}:
{Our Thousand-Year Struggle Over Technology} and {Prosperity}}. Hachette
UK.

\bibitem[\citeproctext]{ref-kemp2025Digital2025Global2025}
Kemp, Simon. 2025. {``Digital 2025: {Global Overview Report}.''}
\emph{DataReportal -- Global Digital Insights}.
https://datareportal.com/reports/digital-2025-global-overview-report;
Data Reportal.

\bibitem[\citeproctext]{ref-kohnMeasurementTrustAutomation2021}
Kohn, Spencer C., Ewart J. de Visser, Eva Wiese, Yi-Ching Lee, and Tyler
H. Shaw. 2021. {``Measurement of {Trust} in {Automation}: {A Narrative
Review} and {Reference Guide}.''} \emph{Frontiers in Psychology} 12
(October). \url{https://doi.org/10.3389/fpsyg.2021.604977}.

\bibitem[\citeproctext]{ref-kongDevelopingValidatingScale2025}
Kong, Siu Cheung, Jinyu Zhu, and Yin Nicole Yang. 2025. {``Developing
and Validating a Scale of Empowerment in Using Artificial Intelligence
for Problem-Solving for Senior Secondary and University Students.''}
\emph{Computers and Education: Artificial Intelligence} 8 (June):
100359. \url{https://doi.org/10.1016/j.caeai.2024.100359}.

\bibitem[\citeproctext]{ref-legaspiSyntheticAgencySense2019}
Legaspi, Roberto, Zhengqi He, and Taro Toyoizumi. 2019. {``Synthetic
Agency: Sense of Agency in Artificial Intelligence.''} \emph{Current
Opinion in Behavioral Sciences}, Artificial {Intelligence}, 29
(October): 84--90. \url{https://doi.org/10.1016/j.cobeha.2019.04.004}.

\bibitem[\citeproctext]{ref-mead1932philosophy}
Mead, George Herbert. 1932. {``The Philosophy of the Present.''}

\bibitem[\citeproctext]{ref-morrowWhenTechnologiesMakes2014}
Morrow, David R. 2014. {``When {Technologies Makes Good People Do Bad
Things}: {Another Argument Against} the {Value-Neutrality} of
{Technologies}.''} \emph{Science and Engineering Ethics} 20 (2):
329--43. \url{https://doi.org/10.1007/s11948-013-9464-1}.

\bibitem[\citeproctext]{ref-moutaWhereAgencyMoving2025}
Mouta, Ana, Ana María Pinto-Llorente, and Eva María Torrecilla-Sánchez.
2025. {``{`{Where} Is {Agency Moving} to?'}: {Exploring} the {Interplay}
Between {AI Technologies} in {Education} and {Human Agency}.''}
\emph{Digital Society} 4 (2): 49.
\url{https://doi.org/10.1007/s44206-025-00203-9}.

\bibitem[\citeproctext]{ref-naeemStepbyStepProcessThematic2023}
Naeem, Muhammad, Wilson Ozuem, Kerry Howell, and Silvia Ranfagni. 2023.
{``A {Step-by-Step Process} of {Thematic Analysis} to {Develop} a
{Conceptual Model} in {Qualitative Research}.''} \emph{International
Journal of Qualitative Methods} 22 (October): 16094069231205789.
\url{https://doi.org/10.1177/16094069231205789}.

\bibitem[\citeproctext]{ref-naikChatGPTAllYou2024}
Naik, Ishita, Dishita Naik, and Nitin Naik. 2024. {``{ChatGPT Is All You
Need}: {Untangling Its Underlying AI Models}, {Architecture}, {Training
Procedure}, {Capabilities}, {Limitations And Applications}.''}
Preprints.
\url{https://doi.org/10.36227/techrxiv.173273427.76836200/v1}.

\bibitem[\citeproctext]{ref-pewFutureHuman2023}
Pew Research Center, Janna Anderson, and Lee Rainie. 2023. {``The
{Future} of {Human Agency},''} February.

\bibitem[\citeproctext]{ref-pinskiAILiteracyUsers2024}
Pinski, Marc, and Alexander Benlian. 2024. {``{AI} Literacy for Users --
{A} Comprehensive Review and Future Research Directions of Learning
Methods, Components, and Effects.''} \emph{Computers in Human Behavior:
Artificial Humans} 2 (1): 100062.
\url{https://doi.org/10.1016/j.chbah.2024.100062}.

\bibitem[\citeproctext]{ref-plotkinaUnearthingAICoaching2024}
Plotkina, Lidia, and Subramaniam Sri Ramalu. 2024. {``Unearthing {AI}
Coaching Chatbots Capabilities for Professional Coaching: A Systematic
Literature Review.''} \emph{Journal of Management Development} 43 (6):
833--48. \url{https://doi.org/10.1108/JMD-06-2024-0182}.

\bibitem[\citeproctext]{ref-PradosdelaEscosura_2022}
Prados de la Escosura, Leandro. 2022. \emph{Human Development and the
Path to Freedom: 1870 to the Present}. New Approaches to Economic and
Social History. Cambridge: Cambridge University Press.

\bibitem[\citeproctext]{ref-selwynDigitalNativeMyth2009}
Selwyn, Neil. 2009. {``The Digital Native -- Myth and Reality.''}
\emph{Aslib Proceedings} 61 (4): 364--79.
\url{https://doi.org/10.1108/00012530910973776}.

\bibitem[\citeproctext]{ref-senDevelopmentFreedom1999}
Sen, Amartya. 1999. \emph{Development as {Freedom}}. Oxford University
Press.

\bibitem[\citeproctext]{ref-silvennoinenStudentsUsageGenAI2025}
Silvennoinen, Minna, Satu Aksovaara, and Merja Alanko-Turunen. 2025.
{``Students' {Usage} of {GenAI} in {Universities} of {Applied Sciences}:
{Experiences} and {Development Needs} for {Guidance} and {Support}.''}
In \emph{38th {Bled eConference}: {Empowering Transformation}: {Shaping
Digital Futures} for {All}: {Conference Proceedings}}, 467--82.
University of Maribor Press.
\url{https://doi.org/10.18690/um.fov.4.2025.29}.

\bibitem[\citeproctext]{ref-sundarRiseMachineAgency2020}
Sundar, S Shyam. 2020. {``Rise of {Machine Agency}: {A Framework} for
{Studying} the {Psychology} of {Human}--{AI Interaction} ({HAII}).''}
\emph{Journal of Computer-Mediated Communication} 25 (1): 74--88.
\url{https://doi.org/10.1093/jcmc/zmz026}.

\bibitem[\citeproctext]{ref-suonpaaStudentsPerceptionsGenerative2024}
Suonpää, Maija, Jutta Heikkilä, and Ana Dimkar. 2024. {``Students'
{Perceptions Of Generative Ai Usage And Risks In A Finnish Higher
Education Institution}.''} In \emph{18th {International Technology},
{Education} and {Development Conference}}, 3071--77. Valencia, Spain.
\url{https://doi.org/10.21125/inted.2024.0825}.

\bibitem[\citeproctext]{ref-surianoStudentInteractionChatGPT2025}
Suriano, Rossella, Alessio Plebe, Alessandro Acciai, and Rosa Angela
Fabio. 2025. {``Student Interaction with {ChatGPT} Can Promote Complex
Critical Thinking Skills.''} \emph{Learning and Instruction} 95
(February): 102011.
\url{https://doi.org/10.1016/j.learninstruc.2024.102011}.

\bibitem[\citeproctext]{ref-swanepoelDoesArtificialIntelligence2021}
Swanepoel, Danielle. 2021. {``Does {Artificial Intelligence Have
Agency}?''} In \emph{The {Mind-Technology Problem}}, edited by Robert W.
Clowes, Klaus Gärtner, and Inês Hipólito, 18:83--104. Cham: Springer
International Publishing.
\url{https://doi.org/10.1007/978-3-030-72644-7_4}.

\bibitem[\citeproctext]{ref-tapalsenseofagencyscale2017}
Tapal, Adam, Ela Oren, and Baruch Eitam. 2017. {``The Sense of Agency
Scale: A Measure of Consciously Perceived Control over One's Mind, Body,
and the Immediate Environment.''} \emph{Frontiers in Psychology} 8
(September): 1552. \url{https://doi.org/10.3389/fpsyg.2017.01552}.

\bibitem[\citeproctext]{ref-winnerCultInnovationIts2018}
Winner, Langdon. 2018. {``The {Cult} of {Innovation}: {Its Myths} and
{Rituals}.''} In \emph{Engineering a {Better Future}: {Interplay}
Between {Engineering}, {Social Sciences}, and {Innovation}}, edited by
Eswaran Subrahmanian, Toluwalogo Odumosu, and Jeffrey Y. Tsao, 61--73.
Cham: Springer International Publishing.
\url{https://doi.org/10.1007/978-3-319-91134-2_8}.

\end{CSLReferences}




\end{document}
